\chapter{Conclusion et perspectives}

Le projet a transformé un squelette UML en un assistant GraphRAG complet. Le métamodèle Arianes est représenté dans Neo4j, les procédures sont découpées et indexées, puis un retriever contrôlé choisit les classes, relations et passages transmis au fournisseur sélectionné. L'application ajoute une mémoire conversationnelle et une séparation des données par utilisateur.

La couche métier constitue un second apport. Elle ne fusionne pas les concepts fonctionnels avec les classes UML, mais les relie en conservant leur provenance. Le pilote autour du poste d'aiguillage montre qu'une réponse peut distinguer une commande fonctionnelle d'une association du métamodèle et citer le passage qui la justifie.

La démarche de validation est volontairement progressive : les 29 fiches ont été extraites et le pilote documentaire global est actif, avec une séparation explicite entre relations confirmées documentairement et relations probables. Ce choix rend visible l'incertitude et prépare la validation experte.

Les perspectives prioritaires sont :

\begin{enumerate}
  \item améliorer la sélection des classes candidates par CamelCase, synonymes et embeddings ;
  \item créer une interface de revue des relations \code{PROBABLE} ;
  \item faire annoter un jeu de questions et de preuves par des experts ;
  \item faire certifier le graphe métier fiche par fiche par des experts ;
  \item mesurer la fidélité des réponses aux citations ;
  \item renforcer le déploiement : HTTPS, sessions, observabilité et base serveur ;
  \item mesurer comparativement la qualité, le coût et la latence de Mistral, Gemini et Groq.
\end{enumerate}

\resultat{L'assistant ne se limite plus à générer du texte : il identifie la classe, la relation et le passage qui justifient la réponse.}
